% === Approach ===
\section{Approach}
\label{sec:approach}

\sys realizes the task of \S\ref{sec:background} as a capability-aware system
organized around retrieval from the memory of resolved incidents. Two commitments
shape it: \emph{match by analogy}, ranking past incidents by similarity to the
current window and reusing their diagnosis; and \emph{adapt to the evidence each
window carries}, so one system spans conditions from a richly instrumented
deployment to a one-line text report. Figure~\ref{fig:arch} shows the pipeline;
throughout, we return the $k=5$ most similar past incidents for each window.

\begin{figure*}[t]
\centering
\includegraphics[width=.9\textwidth]{figures/arise-architecture.pdf}
\caption{The \sys{} instrument. \circled{1}~Each observation window (rich
telemetry or a bare text report) is profiled into eleven binary capability
flags, from which a deterministic controller assembles a cheap-first plan under
a hard per-window budget. \circled{2}~When the fast path is unconvincing
(numeric triage ${<}\,0.90$ or no confident dense match), the escalation gate
schedules the full hybrid retrieval: dense, sparse, and graph retrievers over a
time-aware memory of resolved tickets, fused by reciprocal rank fusion
($\kappa{=}60$, top-20 per retriever); a second gate triggers query
reformulation when $\mathrm{maxconf}(R)<0.5$. \circled{3}~Bounded
evidence-gathering refines the fused ranking, \circled{4}~a calibrated head
issues the triage and novelty decisions, and \circled{5}~a state layer
suppresses duplicate pages across an incident's windows. Dashed arrows are
offline or asynchronous: the knowledge graph is built once per corpus by LLM
extraction over ticket text, and retrievers see only tickets resolved before
the window's timestamp. The gold card summarizes what we \emph{measure} about
the instrument across five training seeds (\S\ref{sec:results-rq1},
\S\ref{sec:results-rq3}).}
\label{fig:arch}
\end{figure*}

\subsection{Capability Profiling and Gated Planning}
\label{sec:approach-controller}

\sys{} first decides \emph{what to do}, then executes; the decision is driven by a
\emph{capability profile}. A lightweight observer maps a window
$w=(t,\mathbf{x},e)$ to a set of binary capability flags
$\mathcal{C}(w)\subseteq\mathcal{F}$, where $\mathcal{F}$ is a fixed vocabulary of
eleven flags (\textsc{Text}, \textsc{NumericFeatures}, \textsc{OrderedLogs},
\textsc{TraceSummary}, \textsc{ClusterEvents}, \textsc{Metrics},
\textsc{GraphEntities}, and others), each set by a deterministic predicate on the
evidence present. Every reasoning step, or \emph{skill}, declares a requirement
set $\mathrm{req}(s)$ and is admissible for $w$ only when
$\mathrm{req}(s)\subseteq\mathcal{C}(w)$. A deterministic controller composes the
admissible skills into an ordered \emph{plan} $P(w)$, choosing a branch from the
window's role in the incident lifecycle: a quiet pre-fault baseline, an active
fault, a recovery phase, or a continuation of an already-paged incident. Because
inadmissible skills are never scheduled, the system degrades gracefully by
construction: a fully instrumented window yields a plan with numeric triage, all
three retrievers, on-demand evidence requests, and an optional verifier, while a
text-only window drives the identical controller to a compact dense-and-graph plan,
with no code or configuration change.

Planning is also cost-aware. Skills are ordered cheapest first, and the expensive
ones sit behind two \emph{gates} that open only under uncertainty.

\noindent\emph{Confidence.} Both gates read a confidence $c\in[0,1]$, the score a
skill attaches to its own output: for the numeric triage classifier its
calibrated ticket-worthiness probability, and for a retriever the calibrated
probability its logistic head assigns to the window given its best match. We
write
\begin{equation}
\mathrm{maxconf}(R)=\max_{r\in\{\text{dense},\,\text{sparse},\,\text{graph}\}} c_r,
\label{eq:maxconf}
\end{equation}
the largest such confidence over the retrievers that have run. We stress that
$\mathrm{maxconf}$ is \emph{not} the fused score of Eq.~\ref{eq:rrf}: reciprocal
rank fusion produces an unnormalized quantity bounded above by
$3/(\kappa+1)\approx0.049$, which would make a comparison against $0.5$
meaningless. The two live on different scales and are never compared. On the real
corpus the realized $\mathrm{maxconf}$ spans $[0.00,0.57]$ with median $0.42$ and
quartiles $0.36$ and $0.48$, so a $0.5$ threshold sits just above the upper
quartile of the distribution it is applied to.

\noindent\emph{The gates.} A fast pass runs the numeric triage classifier and a
dense lookup. The \emph{escalation gate} schedules the costlier sparse and graph
retrievers and the language-model verifier when this cheap path is unconvincing:
when the numeric triage score falls below $0.90$ or the dense retriever returns no
confident match. A text-only window has no numeric signal to read, so this gate
opens on essentially every window of the real corpus, which is the intended
degraded-mode behaviour. The \emph{reformulation gate} then triggers query
rewriting when $\mathrm{maxconf}(R)<0.5$, signalling that no acceptable precedent
was found; if that persists, the window is flagged novel. Both thresholds are the
components' own operating points, fixed in advance rather than tuned on test
data: $0.5$ is the calibrated ticket-worthy/noise boundary shared by every
retrieval head, and $0.90$ is the numeric classifier's high-confidence commit
point. We report their sensitivity in \S\ref{sec:results-rq3}. A
hard per-window budget bounds the worst case, capping language-model tokens
($10^5$), wall-clock time ($90$\,s), monetary cost (\$$0.50$), and total skill
invocations ($12$); an exhausted budget returns the best partial result. Because
almost every skill reads precomputed signals, the only steps that call a language
model at decision time are the gated verifier and reformulator.
Algorithm~\ref{alg:window} summarizes the per-window procedure.

\begin{algorithm}[t]
\caption{Per-window triage in \sys{}}
\label{alg:window}
\small
\begin{algorithmic}[1]
\Require window $w=(t,\mathbf{x},e)$; memory $\mathcal{M}$; budget $B$
\Ensure top-$k$ matches $R$; triage decision; incident id
\State $\mathcal{C}\gets \textsc{Capabilities}(w)$ \Comment{evidence profile}
\If{$\textsc{SuppressedByState}(w)$}
  \State \Return suppressed page with inherited incident id
\EndIf
\State $P\gets \textsc{Plan}(\textsc{Branch}(w),\,\mathcal{C})$ \Comment{admissible, cheap-first}
\State $c\gets \textsc{NumericTriage}(w)$ if $\textsc{Numeric}\!\in\!\mathcal{C}$;\;\; $R_d\gets \textsc{Dense}(w,\mathcal{M}_{<t})$
\If{$c<0.90$ \textbf{ or } $|R_d|<1$} \Comment{escalation gate}
  \State $R\gets \textsc{RRF}\big(\textsc{Dense},\textsc{Sparse},\textsc{Graph};\,\mathcal{M}_{<t}\big)$
  \If{$\mathrm{maxconf}(R)<0.5$} \Comment{Eq.~\ref{eq:maxconf}, not Eq.~\ref{eq:rrf}}
    \State $R\gets \textsc{Reformulate}(w,\mathcal{M}_{<t})$
  \EndIf
  \If{telemetry available \textbf{ and } still uncertain}
    \State $R\gets \textsc{RerankByEvidence}\big(R,\textsc{Gather}(w)\big)$
  \EndIf
\Else
  \State $R\gets R_d$
\EndIf
\State $p\gets \sigma(\boldsymbol{\beta}^{\top}\mathbf{z}+b)$;\;\; decision $\gets [\,p\ge 0.5\,]$
\State id $\gets \textsc{StateUpdate}(w,R_1)$;\;\; \Return $R_{1:k}$, decision, id
\end{algorithmic}
\end{algorithm}

\noindent\textbf{Deterministic control.}  ARISE uses adaptive tool selection but keeps
control deterministic. This preserves reproducibility, bounded cost,
and auditability: each action follows an explicit capability rule and
budget rather than an open-ended model loop. Language-model calls
are therefore limited to gated verification and reformulation steps.

\subsection{Hybrid Memory Retrieval}
\label{sec:approach-retrieval}

The core of \sys{} ranks the eligible memory $\mathcal{M}_{<t}$ against a window
by fusing three complementary measures of similarity, none reliable in isolation.
The \emph{dense} retriever embeds window and ticket text with a fine-tuned
bi-encoder $\phi$ and scores by cosine similarity,
\begin{equation}
s_{\text{dense}}(w,m)=\cos\!\big(\phi(\mathbf{x}_w),\,\phi(\mathbf{d}_m)\big).
\end{equation}
We fine-tune a compact ($384$-dimensional) MiniLM bi-encoder~\cite{wang2020minilm,reimers2019sbert}
with a multiple-negatives ranking loss~\cite{karpukhin2020dpr}: each
window is paired with its gold ticket as a positive against the rest of the batch,
two hard negatives mined by a lexical retriever~\cite{robertson2009bm25}, and one
random negative, teaching it to separate genuinely related incidents from those
that are merely lexically similar. The dense view captures paraphrase and semantic
overlap but blurs exact tokens; the \emph{sparse} retriever, a learned sparse
lexical model~\cite{formal2021splade}, restores that precision by matching the
literal identifiers, error codes, and component names that distinguish otherwise
similar incidents.

The third view is a \emph{knowledge-graph} retriever that captures structural
rather than surface resemblance (\S\ref{sec:approach-kg}). Each retriever proposes
its top $20$ candidates, and we merge the three ranked lists with reciprocal rank
fusion~\cite{cormack2009rrf}. Writing $\mathrm{rank}_r(m)$ for the position of
incident $m$ in retriever $r$'s list, the fused score is
\begin{equation}
s_{\text{fuse}}(m)=\sum_{r\in\{\text{dense},\,\text{sparse},\,\text{graph}\}}
\frac{1}{\kappa+\mathrm{rank}_r(m)},\qquad \kappa=60,
\label{eq:rrf}
\end{equation}
where an incident contributes nothing from a retriever that does not rank it. A final composition step sharpens the head of the list, anchoring on the dense retriever's strongest candidates and promoting those the sparse and graph retrievers independently corroborate. One view's blind spot is thus covered by another's strength, the central reason the fused ranker beats every constituent
alone (\S\ref{sec:results}).

\subsection{Knowledge-Graph Construction and Retrieval}
\label{sec:approach-kg}

We build the knowledge graph once, offline, from the resolved tickets in memory. For each ticket we prompt a lightweight language model
(\texttt{gpt-4o-mini}) to emit a structured fact set: the affected services and
components, any canonical error classes (for example \textsc{OOMKilled} or
\textsc{ConnectionRefused}), a one-sentence root cause, a one-sentence fix with a
coarse fix type, and a short list of observable symptoms. These facts are
materialized as a typed property graph with seven node labels (\textsc{Incident},
\textsc{Service}, \textsc{Component}, \textsc{ErrorClass}, \textsc{RootCause},
\textsc{Fix}, \textsc{Symptom}) and six labelled relations; shared entities are
deduplicated, so two incidents that touch the same component or raise the same
error class are linked through it, encoding the structural neighbourhoods that
recur across an operational history.

At query time the same extraction runs on the window (or a fast rule-based
extractor when no language model is configured), and we score each visible
incident by typed entity overlap. With $A_\theta(\cdot)$ the set of entities of
type $\theta$ for a window or incident, the graph score is a weighted sum of
overlaps,
\begin{equation}
s_{\text{graph}}(w,m)=\sum_{\theta}\lambda_\theta\,
\big|A_\theta(w)\cap A_\theta(m)\big|,
\end{equation}
with weights that reflect each signal's diagnostic value: an exact error class is
weighted most heavily ($\lambda_{\text{error}}=3.0$), followed by the affected
service ($2.0$) and component ($1.5$), while symptoms, which recur across many
unrelated incidents, contribute least ($1.0$). The score is a single graph query
that enforces time-aware visibility directly: an incident is eligible only if
its resolution timestamp precedes the window's.

\subsection{Evidence-Grounded Refinement}
\label{sec:approach-react}

When a window carries rich telemetry but the fused ranking remains uncertain,
\sys{} does what an engineer would: it gathers more evidence before committing.
Following the reason-and-act paradigm~\cite{yao2023react}, it issues a few targeted
requests, each returning one compact piece of evidence rather than a verdict
(cluster events, an extended trace slice, pod-level metrics, or past incidents from
the same fault family), then re-ranks candidates by how well each description agrees
with the gathered evidence. The loop is bounded: repeated identical requests are refused,
and the per-window budget caps the total number of calls. On a text-only window
these requests are never scheduled, so the step is a safe pass-through that returns
the fused ranking unchanged.

\smallskip
\noindent\textbf{A representative trajectory.} On a telemetry-rich active-fault
window (a cart service degraded by a failing Redis dependency), the cheap path
fires with a thin margin: two precedents, a cache outage and a memory leak, score
almost equally. \sys{} escalates to the sparse and graph retrievers, requests pod
events and metrics, observes a cache-pod restart and memory saturation alongside
the latency spike, and re-ranks the corroborated cache-dependency precedent to
rank one, stopping below budget while the state layer folds the window into the
ongoing incident.

\subsection{Triage, Novelty, and Incident-Aware Paging}
\label{sec:approach-state}

The retrieval evidence also drives the paging decision. A lightweight logistic
model combines the available cheap signals, the numeric triage classifier's score
where instrumentation exists and the confidence of the best retrieval match, into
a ticket-worthiness probability $p=\sigma(\boldsymbol{\beta}^{\top}\mathbf{z}+b)$,
and the window is judged ticket-worthy when $p\ge 0.5$; on text-only incidents,
which lack a telemetry classifier, the decision rests on retrieval confidence
alone. A window is flagged \emph{novel}, signalling a failure mode with no
precedent in memory, when no source is confident: the disjunction of a low maximum
retriever confidence, a verifier judgement of novelty when the verifier ran, and an
optional learned novelty score. We treat novelty detection as an auxiliary signal rather than as a fully validated claim on real data, because the text-only real corpus does not provide reliable labels for no-precedent incidents.

Because a single incident unfolds across many windows, paging on each would bury
the on-call engineer. \sys{} keeps a per-service ring buffer of its last twelve
decisions and suppresses a duplicate when a window repeats a recently paged incident
(same top-ranked precedent and fault family within a three-window lookback, no
recovery window between), letting it inherit that incident's identifier. Consecutive
windows of one fault thus fold into a single incident, and the system approaches one
page per incident, not one per window (\S\ref{sec:results}).

The optional language-model verifier, which re-examines a borderline match in natural language,
is itself capability-gated: enabled where calibration shows it helps, disabled where
it hurts; on real engineer-written incidents we find it net-negative and switch it
off. Treating whether an expensive reasoning step helps in a given setting as just
another capability is, we argue, much of what keeps one system robust across very
different evidence regimes.
